\documentclass[11pt]{article}

\usepackage[margin=1in]{geometry}
\usepackage{amsmath,amssymb,amsthm}
\usepackage{graphicx}
\usepackage{hyperref}
\usepackage{xurl}
\usepackage{microtype}
\usepackage{enumitem}
\usepackage{booktabs}
\usepackage{listings}
\usepackage[nameinlink,noabbrev]{cleveref}

\graphicspath{{../figures/}}

\newtheorem{theorem}{Theorem}
\newtheorem{lemma}{Lemma}
\newtheorem{proposition}{Proposition}
\newtheorem{corollary}{Corollary}
\theoremstyle{definition}
\newtheorem{definition}{Definition}
\theoremstyle{remark}
\newtheorem{remark}{Remark}

\newcommand{\TV}{\mathrm{TV}}
\newcommand{\eps}{\varepsilon}

\title{PSC-Q (Addendum):\\Supplementary Evidence and Operator Microstudies}
\author{}
\date{March 23, 2026 (archive v677)}

\begin{document}
\maketitle

\begin{abstract}
This companion note preserves only the supplementary evidence packet for PSC-Q: short operator-facing microstudies, sanity checks, and implementation cautions that were previously shipped as a split supplement.
The main Congestion-EQ paper remains the family theorem/accountant entrypoint and the main PSC-Q paper owns the declared mechanism packet itself; this addendum is optional supporting evidence rather than a prerequisite, and it supplies referee-readable illustrations and operator-facing failure modes without becoming a second verifier, notarization, runtime-conformance, or release-packaging paper.
Figures are omitted in this source-only release; corresponding plots and raw tables remain available on request.\cite{series:artifactpolicy}
\end{abstract}

\paragraph{Scope.}
Optional supporting-evidence addendum to PSC-Q, downstream of Congestion-EQ and not part of the congestion-family entrypoint.\cite{series:congeq,series:pscq}

\paragraph{Imports.}
We import Congestion-EQ as the target witness/accountant contract and PSC-Q as the declared mechanism surface.\cite{series:congeq,series:pscq}
We also import the downstream verifier/notarization/publication stack only as a handoff: replayable verifier bundles belong to W-Congestion-EQ, notarized evaluation logs belong to Eval~1, runtime conformance receipts belong to Operational~B, and public lineage / release packaging belong to Release~A.\cite{series:wcongeq,series:evalnotary,series:opb,series:releaseproperty}

\paragraph{Exports.}
A compact supporting-evidence packet: five short microstudy placeholders, one implementation checklist, and a narrow set of operator-facing cautions for reading PSC-Q deployments.

\paragraph{Boundary.}
This addendum owns only that supporting-evidence packet.
It does \emph{not} own the congestion witness/accountant contract (Congestion-EQ), the declared PSC-Q mechanism packet, the congestion-specific verifier bundle (W-Congestion-EQ), evaluator notarization / retry-prefix logging (Eval~1), runtime-plan conformance receipts (Operational~B), or public lineage / release packaging (Release~A).\cite{series:congeq,series:pscq,series:wcongeq,series:evalnotary,series:opb,series:releaseproperty}

\paragraph{Neighboring-layer stop map.}
\medskip
\noindent\begin{center}
\small
\begin{tabular}{p{0.27\linewidth}p{0.64\linewidth}}
\toprule
Reader need & Earliest owning note \\
\midrule
Recover the omitted-figure placeholders, short microstudy summaries, implementation checklist, and narrow failure-mode notes & This addendum \\
Need the theorem/accountant root for the congestion witness itself & Congestion-EQ \\
Need the declared substitution-padding mechanism packet & PSC-Q \\
Need a replayable congestion-specific checker or verifier bundle & W-Congestion-EQ \\
Need evaluator manifests, attempt records, or notarized retry-prefix evidence & Eval~1 \\
Need runtime-plan conformance receipts for a deployment & Operational~B \\
Need public lineage, release packaging, or release-facing privacy claims & Release~A (with ABOM/OINL where cited downstream) \\
\bottomrule
\end{tabular}
\end{center}

\paragraph{Exported object.}
The exported object is one exact support roster that a referee or operator can quote while keeping PSC-Q itself brief: the five requested archive figure placeholders \texttt{fig\_padding\_worsens\_tv.pdf}, \texttt{fig\_idle\_padding\_partial\_tv.pdf}, \texttt{fig\_service\_bias\_substitution.pdf}, \texttt{fig\_spectral\_dither\_attenuation.pdf}, and \texttt{fig\_utilization\_range\_optimizer.pdf}; one preemptible-idle-filling checklist anchor stating that cover is cancelled or deferred whenever real work arrives; and a narrow failure-mode roster explaining which caution each requested file supports.

\paragraph{Earliest sufficient stop.}
This addendum is complete once a reader can recover the main PSC-Q cautionary examples and cite the checklist/failure-mode notes without reopening the full artifact bundle.
If a deployment needs the mechanism contract itself, a replayable verifier bundle, a notarized evaluator log, runtime conformance receipts, or public release-lineage packaging, the reader should leave this addendum and use PSC-Q, W-Congestion-EQ, Eval~1, Operational~B, and Release~A directly.\cite{series:pscq,series:wcongeq,series:evalnotary,series:opb,series:releaseproperty}

\paragraph{Later terse-paper citation posture.}
For later terse papers under space pressure, default to Congestion-EQ for the witness/accountant root and to PSC-Q for the declared substitution-padding mechanism; cite this addendum only when the exact support packet exported here --- omitted microstudy placeholders, the implementation checklist, or the narrow failure-mode notes --- is itself live.
If the live question is only the public congestion-witness or PSC-Q mechanism claim, stop at those main notes.
If the live question is instead ``where is the supporting microstudy or operator caution that explains that claim?'', widen to this addendum.
If the question becomes replayable verification, runtime conformance, or public release-lineage comparison, leave this addendum and reopen the neighboring owners.
The maintained worked example and paired cutover rule are the general model for that stop / widen / reopen discipline.\cite{series:workedexample,series:pairedterminalcutoverrules}


\begin{table}[t]
\centering
\begin{tabular}{@{}p{0.28\linewidth}p{0.62\linewidth}@{}}
\toprule
Support-packet field & Minimal public content a later terse paper may quote here \\
\midrule
Packet claim & This addendum exports only the PSC-Q support packet: one exact five-file support roster, one implementation checklist anchor, and narrow failure-mode notes. \\
Support roster & \texttt{fig\_padding\_worsens\_tv.pdf}, \texttt{fig\_idle\_padding\_partial\_tv.pdf}, \texttt{fig\_service\_bias\_substitution.pdf}, \texttt{fig\_spectral\_dither\_attenuation.pdf}, and \texttt{fig\_utilization\_range\_optimizer.pdf} are the exact requested archive files in this support packet. \\
Checklist anchor & Cover occupies only otherwise-idle service opportunities; if real work arrives, cover is cancelled or deferred so it never delays real work. \\
Failure-mode roster & \texttt{fig\_padding\_worsens\_tv.pdf} supports the utilization-knee caution; \texttt{fig\_idle\_padding\_partial\_tv.pdf} supports the partial-idle-filling claim; \texttt{fig\_service\_bias\_substitution.pdf} records the shared-queue service-bias pitfall; \texttt{fig\_spectral\_dither\_attenuation.pdf} records why a spectral tier exists; and \texttt{fig\_utilization\_range\_optimizer.pdf} records the utilization-range compilation caution. \\
Wrapper-refresh cutover & Keep this same card only if a successor wrapper changes presentation but keeps that same five-file support roster, the same checklist line, and the same failure-mode roster. \\
How to cite tersely & Quote this card only when the live issue is the supporting microstudy or operator caution itself; otherwise stop at Congestion-EQ or PSC-Q. \\
Do not widen silently & Replayable checkers, notarized evaluator logs, runtime conformance, and release lineage still belong to W-Congestion-EQ, Eval~1, Operational~B, and Release~A. \\
\bottomrule
\end{tabular}
\caption{A minimal public PSC-Q support-packet card. This is the smallest public answer later terse papers can quote when the live issue is the optional supporting microstudy packet itself rather than the congestion witness/accountant root, declared PSC-Q mechanism packet, replayable checker bundle, runtime conformance receipt, or release lineage.}
\label{tab:pscq-addendum-card}
\end{table}

\paragraph{Tiny support-packet lookup drill.}
Suppose a later short paper only needs to justify the sentence ``near the high-utilization knee, additive padding can worsen congestion-TV, while strict idle-filling already helps if cover never delays real work.''
The smallest public support object is exactly this addendum card plus the two requested archive files \texttt{fig\_padding\_worsens\_tv.pdf} and \texttt{fig\_idle\_padding\_partial\_tv.pdf}, together with the checklist rule that cover is cancelled or deferred whenever real work arrives.
If the later sentence also needs the engineering pitfall that shared-queue ``cover'' can bias service, widen only by one more requested file: \texttt{fig\_service\_bias\_substitution.pdf}.
That is enough to show where the cautionary evidence lives and what implementation discipline the support packet assumes, without widening into Congestion-EQ's theorem/accountant packet, PSC-Q's declared mechanism card, or W-Congestion-EQ's replay bundle.

\paragraph{Paired support/mechanism route.}
Sometimes a later terse paper needs both layers at once: one sentence declaring the live PSC-Q packet and one sentence naming the optional microstudy/checklist packet that explains why that declaration is not empty.
The non-redundant route is therefore the reverse of PSC-Q's own paired citation: cite PSC-Q for the 2 ms renewal-style calendar, substitution-only slack harvesting, observer-resolution contract, and $\rho_{\max}=0.8$ utilization boundary, and cite this addendum only for the narrow support sentence that additive padding can worsen congestion-TV near the utilization knee and that cover is therefore cancelled or deferred whenever real work arrives.\cite{series:pscq}
That paired route keeps the support packet honest: this addendum explains one cautionary evidence/checklist packet, but it does not replace PSC-Q's mechanism declaration, Congestion-EQ's witness/accountant root, or W-Congestion-EQ's replayable checker bundle.

\paragraph{Branched family route.}
The full congestion family can now be cited without orphaning this addendum: Congestion-EQ owns the witness/load-envelope/accountant root, PSC-Q owns the mechanism packet, this addendum owns only the optional support sentence about the utilization-knee caution and the ``cancel or defer cover when real work arrives'' checklist rule, and W-Congestion-EQ owns the replayable checker bundle.
That branched route makes the addendum's role explicit: it is a support branch beneath PSC-Q, not a shadow theorem, not a second mechanism note, and not a verifier/release substitute.\cite{series:congeq,series:pscq,series:wcongeq}

\paragraph{Same support-packet / refreshed-wrapper cutover.}
A later terse paper may keep this exact support-packet card across a successor wrapper refresh only when the live PSC-Q caution packet is unchanged: the requested five-file support roster stays the same, the checklist anchor still says that cover is cancelled or deferred whenever real work arrives, and the failure-mode roster still stops at the utilization-knee, service-bias, spectral-tier, and utilization-range cautions named above.
If the successor wrapper swaps any requested file in that five-file roster, changes the idle-filling rule, promotes a new service-bias caveat into the public packet, or materially changes which caution bullets are needed to explain the mechanism claim, then the shortcut fails closed and the reader should reopen this addendum rather than silently reusing the old card.

\paragraph{Tiny successor support drill.}
Suppose a successor release rewrites the public PSC-Q wrapper but still points operators to the same five requested archive files, keeps the identical checklist line that cover never delays real work, and repeats the same utilization-knee / service-bias / spectral-tier / utilization-range caution roster.
Then a later terse paper may honestly keep this same support-packet card and say that only the wrapper changed.
But if the successor packet instead replaces \texttt{fig\_idle\_padding\_partial\_tv.pdf} with a new queue-sharing microstudy, drops \texttt{fig\_utilization\_range\_optimizer.pdf} from the packet, or changes the checklist to permit bounded head-of-line blocking, the support packet itself has changed and the later paper must reopen this addendum.

\section{Supplementary evidence and microstudies (formerly a split supplement)}\label{app:microstudies}
This section preserves the operator-facing microstudies and ``gotchas'' that were previously shipped as a separate PSC-Q supplement in the split archive. The unified archive keeps a single-file rule, so this addendum now carries the salient support packet while PSC-Q itself stays mechanism-facing. Figures are omitted in this source-only release; the corresponding plots and raw tables are available on request.

\subsection{Additive padding can worsen congestion-TV in the high-utilization knee}
The main text emphasizes \emph{substitution} (idle-filling) because it can reduce congestion witnesses without increasing mean load.
In contrast, \emph{additive} padding can increase utilization and move the system closer to the slack ``knee,''
at which point small secret-correlated load differences yield large changes in waiting-time and RTT laws.
Fig.~\ref{fig:worsens} (from the research archive) illustrates a simple regime in which additive padding increases the measured congestion-TV.

\begin{figure}[t]
\centering
\fbox{\begin{minipage}{0.92\linewidth}\centering\small Figure omitted in source-only release. Requested file: \texttt{fig\_padding\_worsens\_tv.pdf}.\end{minipage}}
\caption{Archive microstudy: a ``padding can worsen'' regime when additive padding pushes utilization upward.}
\label{fig:worsens}
\end{figure}

\subsection{Partial idle-filling is often enough}
Even partial substitution (filling only a fraction of idle slots) can materially reduce congestion-TV without increasing mean load.
Fig.~\ref{fig:idle} summarizes an archive sweep supporting this empirical claim.

\begin{figure}[t]
\centering
\fbox{\begin{minipage}{0.92\linewidth}\centering\small Figure omitted in source-only release. Requested file: \texttt{fig\_idle\_padding\_partial\_tv.pdf}.\end{minipage}}
\caption{Archive microstudy: partial idle-filling can already reduce congestion-TV without increasing mean load.}
\label{fig:idle}
\end{figure}

\subsection{Substitution must avoid service bias (engineering pitfall)}
Substitution is only ``free'' if cover is confined to otherwise-idle service opportunities.
If cover work competes with real work (shared queue or shared downstream resource), it can bias service and increase delay.
Fig.~\ref{fig:servicebias} illustrates this pitfall under a simple mixed-queue implementation.

\begin{figure}[t]
\centering
\fbox{\begin{minipage}{0.92\linewidth}\centering\small Figure omitted in source-only release. Requested file: \texttt{fig\_service\_bias\_substitution.pdf}.\end{minipage}}
\caption{Archive microstudy: a service-bias pitfall when ``cover'' is implemented as competing work rather than strict idle-filling.
Implementation takeaway: implement substitution as \emph{preemptible} idle-filling (or as a separate low-priority queue that never delays real work).}
\label{fig:servicebias}
\end{figure}

\paragraph{Implementation checklist (non-exhaustive).}
\begin{itemize}[leftmargin=1.2em]
\item Define at most one service opportunity per slot; do not create extra opportunities when real work is present.
\item Treat cover as cancelable/preemptible: if real work arrives, cancel or defer cover so it never delays real work.
\item If a separate queue is used, enforce strict priority for real work and bound head-of-line blocking (e.g., fixed-size cover frames).
\item Evaluate congestion witnesses across hops (RTT time series, queue delay, loss/ECN) to avoid bottleneck shift.
\end{itemize}
This checklist is therefore support material for reading or operating against the published PSC-Q mechanism packet, not a substitute runtime plan, verifier artifact, evaluator ledger, or release ledger.\cite{series:pscq,series:wcongeq,series:evalnotary,series:opb,series:releaseproperty}

\subsection{Spectral surrogate evidence (why a ``spectral tier'' exists)}
The PSC-Q main text treats a TV-tier contract and flags periodic/spectral tests as a realistic additional tier.
Fig.~\ref{fig:spectral} is an archive ``spectral surrogate'' plot illustrating that declared dither attenuates periodic witnesses,
supporting the motivation for a two-tier view.

\begin{figure}[t]
\centering
\fbox{\begin{minipage}{0.92\linewidth}\centering\small Figure omitted in source-only release. Requested file: \texttt{fig\_spectral\_dither\_attenuation.pdf}.\end{minipage}}
\caption{Archive evidence: dither attenuates periodic witnesses in a spectral surrogate.}
\label{fig:spectral}
\end{figure}

\subsection{Interval-robust compilation against a utilization range}
The ``range compilation'' note in the main text can be made more explicit:
compile $(\delta,\alpha)$ against a declared utilization interval $\rho\in[\rho_{\min},\rho_{\max}]$ (rather than a point estimate),
and size for $\rho_{\max}$ so the deployment remains within budget whenever the monitored regime stays within the published interval.
Fig.~\ref{fig:utilrangeopt} illustrates an archive sweep in which the chosen $(\delta,\alpha)$ tightens as $\rho_{\max}$ increases.

\begin{figure}[t]
\centering
\fbox{\begin{minipage}{0.92\linewidth}\centering\small Figure omitted in source-only release. Requested file: \texttt{fig\_utilization\_range\_optimizer.pdf}.\end{minipage}}
\caption{Archive evidence: interval-robust compilation tightens as the assumed upper utilization bound increases.}
\label{fig:utilrangeopt}
\end{figure}

\begin{thebibliography}{99}

\bibitem{series:pscq}
Anonymous.
\newblock PSC-Q: Public Service Calendars and Slack-Harvesting Substitution Padding for Congestion-Private Anonymous Overlays.
\newblock Congestion series Paper~2, 2026. (This archive.)

\bibitem{series:congeq}
Anonymous.
\newblock Congestion-EQ: A Congestion Witness Contract for Anonymous Overlays.
\newblock Congestion series Paper~1, 2026. (This archive.)

\bibitem{series:workedexample}
Anonymous.
\newblock Worked Example: Instantiating a Tiered Reader-Privacy Receipt.
\newblock Synthesis~17, The-One-True-Archive (2026).

\bibitem{series:pairedterminalcutoverrules}
Anonymous.
\newblock Paired Terminal Cutover Rules for Stable Review Spines: When Later Terse Papers Must Stop, Widen, or Reopen.
\newblock Series synthesis note (Synthesis~75), 2026. (This archive.)

\bibitem{series:artifactpolicy}
Anonymous.
\newblock Artifact policy and supporting materials for the anonymity/AnonDHT series.
\newblock 2026.

\bibitem{series:wcongeq}
Anonymous.
\newblock \emph{W-Congestion-EQ: Verifiable Congestion Budgets --- Congestion-Specific Verifier Bundles for Anonymous Overlays}.
\newblock Congestion series Paper~3, 2026. (This archive.)

\bibitem{series:evalnotary}
Anonymous.
\newblock \emph{Notarized Anonymity Certificates Under Retries: Grinding-Resistant Evaluation via Transparency Logs, Public Randomness, and $\delta$-Spending}.
\newblock Evaluation series Paper~1, 2026. (This archive.)

\bibitem{series:opb}
Anonymous.
\newblock \emph{Runtime Conformance Receipts for Anonymous DHTs: Plan Pinning, PTL Checks, and Epoch Trace Audits}.
\newblock Operational series Paper~B, 2026. (This archive.)

\bibitem{series:releaseproperty}
Anonymous.
\newblock \emph{Privacy as a Release Property: Ledger Bundles and Receipts for Anonymous DHT Deployments}.
\newblock Release-and-destination Paper~A, 2026. (This archive.)
\end{thebibliography}

\end{document}
